\section{Contextual Bandits}
\label{sec:exp3_results}

Experiment~3 replaces Q-learning with contextual bandits. LinUCB uses optimism-based exploration while Contextual Thompson Sampling uses posterior sampling, creating a natural contrast for how exploration style affects outcomes. Each receives a tailored factorial design to avoid wasting design cells on factors with no effect in the other algorithm.

\subsection{LinUCB}

Experiment~IIIa deploys LinUCB agents under the affiliated valuation model. The design is a $3 \times 2^7 = 384$ mixed-level factorial with 9 factors (7 binary plus the three-level affiliation parameter $\eta$), including the regularisation parameter $\lambda$ and memory decay, both of which strongly affect LinUCB performance. The design is replicated twice for \ExpThreeANobs{} observations.

\begin{figure}[H]
  \centering
  \includegraphics[width=\textwidth]{figures/e3a_trace}
  \caption{Representative bid function for a single trial of Experiment~IIIa (LinUCB, first-price auction, $\eta = 0.5$, 2~bidders, 5{,}000 rounds). Bids versus signals in the final 1{,}000 rounds, overlaid with the BNE bid function.}
  \label{fig:e3a_trace}
\end{figure}

\subsubsection{Revenue}

Number of bidders again dominates ($|t| = \ExpThreeARevFactorTOne$), with the shift from two to four bidders increasing revenue by \ExpThreeARevTopPctEffect\% of the grand mean. The next largest effects are \ExpThreeARevFactorTwo{} ($|t| = \ExpThreeARevFactorTTwo$) and \ExpThreeARevFactorThree{} ($|t| = \ExpThreeARevFactorTThree$). Table~\ref{tab:exp3a_ranked_rev} ranks all significant effects by absolute $t$-statistic. First-price auctions yield lower final revenue ($|t| = \ExpThreeARevAuctionT$), as the main effects plot (Figure~\ref{fig:e3a_main_rev}) confirms.

\begin{table}[H]
\centering
\caption{Experiment 3a: Significant effects for average revenue ($p < 0.05$), ranked by $|t|$.}
\label{tab:exp3a_ranked_rev}
\begin{tabular}{lrrl}
\toprule
\textbf{Effect} & \textbf{Coeff.} & \textbf{$|t|$} & \textbf{Direction} \\
\midrule
Budget multiplier & 2006.6019 & 39.44 & + \\
Bidder objective & -1428.0108 & 28.07 & $-$ \\
Bidder objective $\times$ Budget multiplier & -1387.4815 & 27.27 & $-$ \\
Number of bidders & 1158.2655 & 22.76 & + \\
Bidder objective $\times$ Number of bidders & -629.6958 & 12.38 & $-$ \\
Number of bidders $\times$ Budget multiplier & 448.9495 & 8.82 & + \\
Value dispersion ($\sigma$) & 356.7849 & 7.01 & + \\
Auction format & 184.0336 & 3.62 & + \\
Auction format $\times$ Bidder objective & 182.7252 & 3.59 & + \\
Auction format $\times$ Budget multiplier & 140.8645 & 2.77 & + \\
\bottomrule
\end{tabular}
\end{table}


\begin{figure}[H]
  \centering
  \includegraphics[width=0.7\textwidth]{figures/e3a_main_rev}
  \caption{Experiment~IIIa: Main effects plot for average revenue under LinUCB.}
  \label{fig:e3a_main_rev}
\end{figure}

Figure~\ref{fig:e3a_int_rev} displays the top interaction effects.

\begin{figure}[H]
  \centering
  \includegraphics[width=0.7\textwidth]{figures/e3a_int_rev}
  \caption{Experiment~IIIa: Interaction plot for average revenue under LinUCB. Non-parallel lines indicate factor interdependencies.}
  \label{fig:e3a_int_rev}
\end{figure}

\subsubsection{Price Volatility}

Table~\ref{tab:exp3a_ranked_vol} and the main effects plot (Figure~\ref{fig:e3a_main_vol}) show the factor rankings for price volatility under LinUCB.

\begin{table}[H]
\centering
\caption{Experiment 3a: Significant effects for price volatility ($p < 0.05$), ranked by $|t|$.}
\label{tab:exp3a_ranked_vol}
\begin{tabular}{lrrl}
\toprule
\textbf{Effect} & \textbf{Coeff.} & \textbf{$|t|$} & \textbf{Direction} \\
\midrule
Budget multiplier & 1.2878 & 3.26 & + \\
Bidder objective & -1.1490 & 2.91 & $-$ \\
Bidder objective $\times$ Budget multiplier & -1.1266 & 2.85 & $-$ \\
Auction format & -0.8413 & 2.13 & $-$ \\
Auction format $\times$ Bidder objective & 0.7996 & 2.02 & + \\
Auction format $\times$ Value dispersion ($\sigma$) & -0.7853 & 1.99 & $-$ \\
\bottomrule
\end{tabular}
\end{table}


\begin{figure}[H]
  \centering
  \includegraphics[width=0.7\textwidth]{figures/e3a_main_vol}
  \caption{Experiment~IIIa: Main effects plot for price volatility under LinUCB.}
  \label{fig:e3a_main_vol}
\end{figure}

Model adequacy diagnostics (Appendix~\ref{sec:appendix_robustness}) confirm that the linear model captures the dominant effects.\footnote{$R^2$ values reach \ExpThreeARsqMax{} across \ExpThreeANobs{} observations. Of the \ExpThreeANRawSig{} effects significant at $p < 0.05$, \ExpThreeANBHSig{} (\ExpThreeABHPct\%) survive Benjamini--Hochberg correction.}

\subsection{Contextual Thompson Sampling}

Experiment~IIIb deploys Contextual Thompson Sampling agents under the same affiliated valuation model. The design is a $3 \times 2^5 = 96$ mixed-level factorial with 7 factors (5 binary plus $\eta$). Preliminary screening confirmed that $\lambda$ and memory decay have no detectable effect on Thompson Sampling outcomes, so both are excluded to avoid wasting design cells on inert factors. The design is replicated twice for \ExpThreeBNobs{} observations.

\begin{figure}[H]
  \centering
  \includegraphics[width=\textwidth]{figures/e3b_trace}
  \caption{Representative bid function for a single trial of Experiment~IIIb (Thompson Sampling, first-price auction, $\eta = 0.5$, 2~bidders, 5{,}000 rounds).}
  \label{fig:e3b_trace}
\end{figure}

\subsubsection{Revenue}

Under Thompson Sampling, \ExpThreeBRevFactorOne{} dominates ($|t| = \ExpThreeBRevFactorTOne$), with the effect amounting to \ExpThreeBRevTopPctEffect\% of the grand mean, followed by \ExpThreeBRevFactorTwo{} ($|t| = \ExpThreeBRevFactorTTwo$). Table~\ref{tab:exp3b_ranked_rev} ranks all significant effects.

\begin{table}[H]
\centering
\caption{Experiment 3b: Significant effects for average revenue ($p < 0.05$), ranked by $|t|$.}
\label{tab:exp3b_ranked_rev}
\begin{tabular}{lrrl}
\toprule
\textbf{Effect} & \textbf{Coeff.} & \textbf{$|t|$} & \textbf{Direction} \\
\midrule
Budget multiplier & 1316.1821 & 31.68 & + \\
Number of bidders & 667.7779 & 16.07 & + \\
Auction format & 458.3301 & 11.03 & + \\
Auction format $\times$ Budget multiplier & 393.6340 & 9.47 & + \\
Value dispersion ($\sigma$) & 360.6121 & 8.68 & + \\
Auction format $\times$ Value dispersion ($\sigma$) & 233.6535 & 5.62 & + \\
Budget multiplier $\times$ Value dispersion ($\sigma$) & 175.0143 & 4.21 & + \\
Number of bidders $\times$ Value dispersion ($\sigma$) & 146.9008 & 3.54 & + \\
\bottomrule
\end{tabular}
\end{table}


\begin{figure}[H]
  \centering
  \includegraphics[width=0.7\textwidth]{figures/e3b_main_rev}
  \caption{Experiment~IIIb: Main effects plot for average revenue under Thompson Sampling.}
  \label{fig:e3b_main_rev}
\end{figure}

\begin{figure}[H]
  \centering
  \includegraphics[width=0.7\textwidth]{figures/e3b_int_rev}
  \caption{Experiment~IIIb: Interaction plot for average revenue under Thompson Sampling.}
  \label{fig:e3b_int_rev}
\end{figure}

\subsubsection{Price Volatility}

Table~\ref{tab:exp3b_ranked_vol} and the main effects plot (Figure~\ref{fig:e3b_main_vol}) show the factor rankings for price volatility under Thompson Sampling.

\begin{table}[H]
\centering
\caption{Experiment 3b: Significant effects for price volatility ($p < 0.05$), ranked by $|t|$.}
\label{tab:exp3b_ranked_vol}
\begin{tabular}{lrrl}
\toprule
\textbf{Effect} & \textbf{Coeff.} & \textbf{$|t|$} & \textbf{Direction} \\
\midrule
Budget multiplier & 0.3873 & 54.37 & + \\
Auction format & -0.1585 & 22.25 & $-$ \\
Auction format $\times$ Budget multiplier & 0.1233 & 17.31 & + \\
Auction format $\times$ Number of bidders & 0.0410 & 5.76 & + \\
\bottomrule
\end{tabular}
\end{table}


\begin{figure}[H]
  \centering
  \includegraphics[width=0.7\textwidth]{figures/e3b_main_vol}
  \caption{Experiment~IIIb: Main effects plot for price volatility under Thompson Sampling.}
  \label{fig:e3b_main_vol}
\end{figure}

Model adequacy diagnostics confirm that the linear model captures the dominant effects.\footnote{$R^2$ values reach \ExpThreeBRsqMax{} across \ExpThreeBNobs{} observations. Of the \ExpThreeBNRawSig{} effects significant at $p < 0.05$, \ExpThreeBNBHSig{} (\ExpThreeBBHPct\%) survive Benjamini--Hochberg correction.}

\subsection{Summary}

Replacing Q-learning with contextual bandits does not improve seller welfare. Richer algorithm design does not translate into higher revenue for the auctioneer. LinUCB and Thompson Sampling respond to different design levers, validating the decision to analyse them separately. Higher exploration intensity reduces revenue under both algorithms, but the magnitudes and moderating factors differ. Global sensitivity analysis confirms these hierarchies. Under LinUCB, the number of bidders accounts for over half of revenue variance ($S_T = 0.55$) and nearly all winner entropy variance ($S_T = 0.75$). Under Thompson Sampling, factor importance is more diffuse, with no single factor exceeding $S_T = 0.31$ for revenue. Across both algorithms, tuning parameters are consistently negligible ($S_T < 0.04$ for revenue), reinforcing the conclusion that market structure dominates algorithm configuration (Appendix~\ref{sec:appendix_sensitivity}). Experiments~1--3 share a common feature, namely that agents bid directly from value estimates without external constraints. Section~\ref{sec:exp4_results} tests whether the patterns documented above persist when agents face budget constraints, the dominant paradigm in modern advertising exchanges.
